% !TeX root = ../main.tex
\begin{frame}{04 结论与展望——挑战与展望}
  \begin{SlideCanvas}
    \BodyAt{55}{84}{850}{联合编程为器件设计提供方法基础，进一步应用需衔接径向部署、真实材料与医学功能评价。}
    \PointAt{55}{148}{235}{挑战}
    \PointAt{337}{148}{275}{发展方向}
    \PointAt{650}{148}{250}{本方案的落点}
    \MainRuleAt{186}
    \PointAt{55}{209}{245}{多项响应相互耦合}
    \DetailAt{337}{209}{275}{全过程的多目标设计，报告误差及可行范围}
    \DetailAt{650}{209}{250}{长度与承载共同评价}
    \RuleAt{264}
    \PointAt{55}{287}{245}{器件工况的迁移}
    \DetailAt{337}{287}{275}{轴向拉伸向径向部署扩展，必要时直接器件优化}
    \DetailAt{650}{287}{250}{补充径向加载与端部检查}
    \RuleAt{342}
    \PointAt{55}{365}{245}{材料与制造受约束}
    \DetailAt{337}{365}{275}{实测几何与相变参数反馈，纳入缺陷与长期载荷}
    \DetailAt{650}{365}{250}{局部应变、体温功能及疲劳}
    \RuleAt{420}
    \PointAt{55}{438}{245}{医学功能评价}
    \DetailAt{337}{438}{275}{接触、血流及生物学评价}
    \DetailAt{650}{438}{250}{接触改善与生物反应验证}
    \Sources{制造与材料：\RefVella；\RefYasmin。机械—医学机制：\href{https://doi.org/10.1016/j.jscai.2023.100971}{SCAI共识，2023}；接触评价：\href{https://doi.org/10.1007/s11665-014-0913-z}{Borghi等，2014}。}
  \end{SlideCanvas}
  \BICConclusionBand{器件应用需要\BICBandEmph{几何、材料与医学评价}共同支撑}
  \note{趋势依据代表性研究综合归纳，不声称现有方法普遍缺少某项能力。Vellaparambil2024说明实测杆厚需要反馈有限元；Yasmin2026说明打印的NiTi几何可以成形，但其成形件Af高于37摄氏度，体温下的超弹性受到限制，需后处理和材料复测；故几何参数和材料状态必须共同验证。单胞与器件差异还包括曲率、端部、接触和加载方向。再狭窄或血栓结局需要相应生物和临床研究，不能由泊松比或应力云图直接推出。本方案先建立长度与径向承载的可比较证据，再决定是否扩展至接触和血流。}
\end{frame}
